\author{Łukasz Ziobroń}
\title[Crisis]{(R)evolution of C++}
\institute{lukasz@ziobron.net \and \url{http://ziobron.net}}
\date{code::dive, 2016-11-15}
\subject{Computer Science}
\subtitle{The Ultimate Hitchhiker's Guide to the C++}
%\logo{}

\begin{frame}
\titlepage
\end{frame}

\begin{frame}
  \frametitle{(R)evolution of JS}
  \centering 
  \includegraphics[height=0.8\paperheight]{writing-device-drivers-with-js}
\end{frame}

\begin{frame}
  \frametitle{About the author}
  \includegraphics[height=50pt]{html4-logo}
  \includegraphics[height=40pt]{php-logo}
  \includegraphics[height=60pt]{c++-logo}
\end{frame}


\slide{Key messages}{
  \begin{enumerate}
    \item C++ had a clear aim, which made it popular: to organize code better without the loss of efficiency
    \item C++ is even more popular now, because of new standards: C++11 and C++14.
    \item C++ will be one of the most popular programming languages in future so it's worth to invest your time and/or money to learn it.
  \end{enumerate}
}

%\slide{Agenda}
%{
%\tableofcontents
%}

% note that this way you can easily comment-out compleated sections, when working on the document.
% this signifficantly speeds up work. oh - and use 'make fast' target for compiltion - it will run latex just once.
% it may result in incorrect indexes and numbers missmatch, but for a preview is just fine.
%\input{autonomous_cars.tex}
%\input{image_processing.tex}
%\input{machine_learning.tex}
%\input{tools.tex}


